\begin{abstract}
Optimization-based bound tightening can reduce a search tree while
increasing total solving time. This report develops finite checks that
separate possible tightening benefit from computational cost. Feasible
points bound the benefit of a frozen relaxation round. A finite pool of
points feasible in the relaxation rebuilt on its own coordinate hull
protects an inner box through every subsequent round of the same monotone
relaxation family, with explicit cutoff and domain reuse conditions.
A uniform matrix bound yields a separate residual certificate for
remaining endpoint motion. Exact parametric LP regions describe cutoff
sensitivity across active-set changes. A feasible-repair theorem transfers
uniform nonlinear constraint residual and objective-error bounds into
contraction under a quantitative growth condition. Explicit examples show
an exact rate and the failure of extrapolation from observed ratios.
These results use established duality, order, and contraction
arguments, with rational reference checks and explicit scope. A limited
SCIP propagator implements numerical bound validation, current-round
witness filtering, pilot stopping, and event-triggered reconsideration.
The report distinguishes its comparative experiment from implementation
of the stronger future-round certificate machinery. In a prospective
120-run comparison, native, fixed, and adaptive policies each solve the
same 14 of 40 model--seed pairs. Adaptive uses 20.9\% fewer directional
LPs than fixed but 18.4\% more measured sidecar time, with no additional
solves. Native SCIP remains the default. An enhancement
budget controls optional work on the modified search trajectory; it does
not guarantee a runtime bound relative to an independent baseline.
Certificate validity, cheap certificate discovery, and overall solver
performance remain distinct claims.
\end{abstract}
